\documentclass[../main.tex]{subfiles}
\begin{document}
%==========================================TIÊU ĐỀ TẬP========================================
\begin{table}[h]
  \begin{adjustwidth}{-1cm}{}
    \begin{tabular}{>{\centering\arraybackslash}p{6cm}|>{\raggedright\arraybackslash}p{14cm}}
      \multirow{5}{6cm}
      % Cột bên trái: ÉPISODE và số tập
      {\\[-40pt]
        \flaregothic\fontsize{20pt}{20pt}\selectfont \centering
        \color{eptype}HISTOIRE\color{black}\\
        \gangofthree\fontsize{72pt}{72pt}\selectfont
        \color{epnum}2
        \\[18pt]
        % \flaregothic\fontsize{14pt}{14pt}\selectfont
        % \color{epattr}BIENVENUE, \uline{\color{Banpire} Mel} !
        % \flaregothic\fontsize{18pt}{18pt}\selectfont
        % \color{epattr}ÉPISODE PERDU
        \color{black}
      }
       &
      % Dòng thứ nhất: tên tập tiếng Nhật
      {\fotskip\fontsize{18pt}{18pt}\selectfont ようこそ、\ruby{火}{ひ}\ruby{川}{かわ}\ruby{家}{け}へ！}
      \\[6pt]
      % Dòng thứ hai: tên tập tiếng Anh
       & {\cafeteria\fontsize{18pt}{18pt}\selectfont Welcome to the Hikawa Family!}        \\[4pt]
      % Dòng thứ ba: tên tập tiếng Việt
       & {\vnmsans\fontsize{18pt}{18pt}\selectfont Chào mừng đến với nhà Hikawa!}          \\[8pt]
      % Dòng thứ tư: Nhân vật xuất hiện
       & {\brian{\fontsize{14pt}{14pt}\selectfont Track used: Lunar Zoo Year - Grasswalk}}
    \end{tabular}
  \end{adjustwidth}
\end{table}
%===========================================NỘI DUNG==========================================
\color{black}

\begin{center}
  \uline{Hiện tại.}
\end{center}

\dots Nhưng chỉ có điều, cuộc hành trình của hai mươi cô gái đến nhà của Hikawa Kuon không được suôn sẻ lắm. Bởi lẽ ngay khi cánh cửa vừa bật mở, tất cả mọi người (ngoại trừ A-chan và Sora) lao vào trong với một tốc độ đáng kinh ngạc. Cảnh chen lấn hỗn loạn kết thúc bằng việc gần hai mươi cô gái ngã sõng soài xuống sàn gỗ. Tiếng ``rầm, rầm'' vang lên liên tục, tạo nên một màn hỗn độn đáng nhớ.

\begin{dialogue}
  \speak{Mel} Nào mọi người, đừng có chen lấn xô đẩy nữa- Á!
  \speak{Subaru} Ái da, chân chị!
  \speak{Korone} Cái đó là đuôi tôi mừ\~{}
  \speak{Shion} Đau, tay của em!
\end{dialogue}

Trong số họ, Noel và Roboco xui xẻo nhất khi ngã thẳng vào một chậu nước lau nhà bẩn được đặt cạnh tường, khiến nước đổ lênh láng khắp sàn.

``\textbf{Ôi giời ơi, anh vừa mới lau cái sàn xong đó, hai cái đứa PON này!}'' một giọng nói quen thuộc, pha chút bất lực lẫn ngán ngẩm vang lên từ cuối phòng. Nàng Hiệp sĩ và nàng Người máy không khó để nhận ra đó là ai. Cả hai chỉ biết ngượng ngùng cười trừ xin lỗi.

Trong khi đó, nàng ``Tăng ca'' và nàng Thần tượng bước vào phòng với phong thái hoàn toàn trái ngược. Cả hai cúi đầu chào hỏi rất lịch sự, cố gắng làm dịu đi bầu không khí.

\begin{dialogue}
  \speak{A-chan} Tôi xin lỗi vì sự vô kỷ luật của mọi người\dots
  \speak{Sora} Làm phiền cả nhà anh rồi, Kuon-san! Úi chà\dots
\end{dialogue}

Sora vấp phải chân của một trong những người đang ngã ngựa dưới sàn, nhưng may mắn thay cô chống tay được vào tường. Trong khi đó,  Haato và Pekora chạy đến kéo hai ``đứa PON'' dậy. Họ không chần chừ, lấy ngay cây chổi dựng trên tường để lau dọn chỗ nước bẩn vừa làm đổ.

Khi mọi người lấy lại trật tự và nhìn xung quanh, tất cả đều bất ngờ với căn phòng mà họ vừa bước vào. Trái ngược hoàn toàn với những căn phòng họ thường ở trong các buổi quay, nơi đây lại nhỏ bé và giản dị đến khó tin.

Matsuri, Pekora và Shion lên tiếng đồng thanh hỏi trong sự ngỡ ngàng, khiến chủ nhà gãi đầu khó xử:

\begin{dialogue}
  \speak{Matsuri} Nhà của anh\dots
  \speak{Pekora} \dots tại đây\dots
  \speak{Shion} \dots chỉ như vậy thôi sao?
  \speak{Kuon} Ờm\dots\ có lẽ mấy đứa sẽ không tin, nhưng\dots\ ừ. Nhà anh chỉ có như thế này thôi.
\end{dialogue}

Căn phòng này chỉ rộng bằng một nửa (\(\frac{1}{2}\)), thậm chí là một phần ba (\(\frac{1}{3}\)). Nói một cách dễ hiểu, nó khá nhỏ, chỉ ``rộng chưa tới năm mươi mét vuông (\(50m^2\))'' (theo lời của Kuon).

Đó là một căn phòng hình chữ L, có sơ đồ như dưới đây:

\img[\textwidth]{kuon_f5}

Dụng cụ trong nhà không nhiều: một chiếc vô tuyến khoảng 35 inch, một bộ set-top box chạy hệ điều hành A*dro*d, một bộ dàn nhạc với hệ thống loa bình dân (nhưng không quá cổ lỗ sĩ), vài chiếc tủ đựng đồ quần áo, và một tủ kính chứa đồ lưu niệm. Trong phòng còn có một chiếc bàn gỗ, một chiếc đệm vừa đủ cho bốn người nằm, cùng một số vật dụng cơ bản như quạt cây (công tắc đã hỏng, phải dùng tua vít để điều chỉnh), bộ phát mạng, và điều hòa một chiều.

Nàng Phù thủy buột miệng: ``Nhà anh trông có vẻ nghèo với chật chội nhỉ.'' Cậu chủ nhà nở một nụ cười chát chúa, rồi đáp lại bằng giọng trầm ngâm: ``Em không nhất thiết phải nói thẳng toẹt như thế đâu\dots''

Ngay lập tức, một làn sóng phản đối nổi lên bênh vực cậu Tổng từ những người khác.

\begin{dialogue}
  \speak{A-chan} Cẩn thận lời nói của mình vào chứ, Murasaki-san.
  \speak{Choco} Em không cần phải khoáy sâu vào nỗi đau của Kuon-sama như vậy đâu.
  \speak{Rushia} Đ-Đúng đó! Dù đồ đạc hơi ít thật, nhưng nhà anh ấy cũng đâu thiếu thốn đến mức thế!
  \speak{Okayu} Với lại, nhà bé cũng là một lợi thế mừ\~{}
  \speak{Subaru} Nhà bé thì càng ấm cúng hơn chứ, phải không!?
\end{dialogue}

Cậu Tổng thở dài, đứng chống hông lắc đầu bất lực: ``\textit{Mấy đứa đang an ủi anh hay đang trù ẻo anh đấy\dots}''

Bất chợt, nàng Quỷ sừng hướng mặt ra cửa sổ, rồi vội vàng kéo rèm và chỉ tay ra ngoài: ``Mọi người ơi, nhìn ra ngoài đi!''

Cả phòng hướng ánh mắt về phía cửa sổ. Bên ngoài, quang cảnh khu phố hiện lên rõ nét: những xe máy nườm nượp chạy trên đường, những tiếng còi xe bíp bíp inh ỏi trong giờ tan tầm, những dãy nhà chen chúc nằm sát nhau, xen lẫn với các cửa hàng tạp hóa, nhà hàng nhỏ, và đôi ba hiệu sách cũ kỹ. Ở phía xa xa là những tòa nhà cao tầng thỉnh thoảng lọt vào tầm mắt.

Hai con đường chính bị ngăn cách bởi một dòng sông đen ngòm, đầy rác, thoang thoảng bốc lên một mùi hôi thối của nước thải, khiến ai đi qua trông thấy cũng phải nhíu mày khó chịu.

Nhưng với các cô gái, đây là lần đầu tiên họ được trông thấy một thành phố đông đúc như vậy. Vì nhà của cậu khá chật, các cô gái phải chen nhau nhìn từ bốn ô cửa sổ rộng và một ô cửa sổ hẹp.

\begin{dialogue}
  \speak{Aki} Khu nhà của anh trông đông vui thật.
  \speak{Pekora} Trông còn chật ních hơn cả khu phố ta sống nữa, peko.
  \speak{Mio} Bởi vì khu ta làm gì có ai ngoài chúng ta đâu.
  \speak{Subaru} Nhiều xe máy quá!
  \speak{Roboco} Không khí ở đây cũng ngột ngạt hơn hẳn\dots
  \speak{Rushia} Vâng\dots\ Em thấy cũng hơi khó thở hơn\dots
\end{dialogue}

Cậu con trai chẹp miệng, khoanh tay trước ngực và nhún vai như muốn buông xuôi: ``Rất tiếc cho mấy đứa, nhưng trung tâm thủ đô Kawachi này là một trong những thành phố ô nhiễm nhất thế giới đấy. Việc inh ỏi còi xe với tắc đường và ô nhiễm không khí là chuyện thường ở huyện tại đây.''

Tất cả mọi người đều sững sờ trước thông tin này, rồi cùng nhìn cậu với những cặp mắt ái ngại.

``Sao anh nói như thể đó là điều hiển nhiên vậy\dots'' nàng Sói đen lên tiếng, đổ mồ hôi hột.

Cậu nhún vai đáp, như thể đó là lẽ dĩ nhiên: ``Sống lâu thì quen thôi mà em. Ăn bẩn sống lâu, ăn cứt trâu bất tử. Hoặc ít nhất họ nói như vậy.''

Câu nói đùa(?) thô thiển của cậu lại càng khiến cho họ lạnh sống lưng hơn. ``S-sao anh có thể sống được ở đây vậy nhỉ\dots''

Cậu thở dài, rồi tiếp tục giải thích, chỉ tay về phía trước: ``Với cả, chúng ta đang ở trên tầng năm của ngôi nhà này đấy. Và đây cũng là nơi anh sống luôn.''

Những ánh mắt ngơ ngác quay lại nhìn cậu. Mỗi người như mang trên đầu một dấu hỏi to đùng.

``Mà, nhà anh có bao nhiêu tầng vậy?'' nàng Pháp sư tò mò hỏi. Không để mọi người thắc mắc quá lâu, cậu bắt đầu giải thích: ``Gọi là nhà cũng\dots\ không hẳn, nhưng tòa này có sáu tầng đấy. Mỗi tầng lại chia thành hai bên.''

Nàng Cáo trắng nghiêng đầu thắc mắc: ``Nhà anh là dạng chung cư à?'' Kuon khẽ lắc đầu: ``Cũng chẳng phải. Toàn người nhà ở đây hết thôi. Trừ hai tầng dưới cùng. Và tầng sáu nữa.''

``Sao lại thế?'' nàng Cổ động tỏ vẻ khó hiểu, ``Tại sao lại là hai tầng dưới và cả tầng sáu ạ?''

Cậu tiếp tục giải thích: ``Hai tầng dưới nhà anh cho bên dịch vụ viễn thông. Tầng sáu nhà anh là bếp với phòng thờ rồi, ai lại ở trên đó?''

Tất cả mọi người đều gật gù với lời của cậu. Đến lượt Sora đưa lên ý kiến: ``Tại sao nhà này lại chia thành hai bên vậy?''

Cậu khẽ thở dài, ánh mắt trở nên xa xăm hơn, như thể đây là câu hỏi mà cậu muốn tránh né. ``\dots Chuyện nội bộ trong gia đình anh. Không tiện kể đâu.''

Nàng Dơi lên tiếng, kéo theo đó là mọi người đoán già đoán non:

\begin{dialogue}
  \speak{Mel} Kiểu như, tranh chấp tài sản ư?
  \speak{Fubuki} Nội chiến gia đình?
  \speak{Noel} Chiến tranh giành ngôi vị?
\end{dialogue}

Cậu lắc đầu, đưa tay lau trán dang đổ mồ hôi hột: ``Mấy đứa đang nghĩ cái quái gì thế? Mà này, chuyện nhà người ta chúi mũi vào làm cái gì!''

Tất cả phá lên cười, phá tan bầu không khí ngượng ngập. Nhưng bên cạnh tiếng cười, đôi mắt của vài người vẫn không giấu được sự tò mò\dots

Tiếp tục với câu hỏi về nhà của Kuon, nàng Y tá nghiêng đầu hỏi: ``Anh nói rằng nhà anh có sáu tầng, Kuon-sama, với hai tầng dưới cho thuê và tầng sáu là nhà bếp và phòng thờ rồi, vậy các tầng còn lại thì sao?''

Câu hỏi này khiến các cô gái lập tức chú ý, ánh mắt đầy tò mò đổ dồn về phía cậu. Kuon đáp lại một cách điềm nhiên: ``Tầng ba hiện đang bỏ không, tầng bốn bên trái là phòng của ông bà nội anh, còn bên phải là phòng khách của bác anh\dots''

``Còn tầng năm thì sao?'' Aqua hỏi tiếp, mắt mở to chờ đợi. Cậu nhún vai, trả lời ngắn gọn: ``Bên phải tầng này chỉ là phòng ngủ của bác anh thôi.''

Aki nhẹ nhàng lên tiếng: ``Vậy nhà anh chỉ có mỗi chỗ này thôi à?'' Cậu gật đầu, tỏ vẻ hơi đượm buồn: ``Ừ, tất cả sinh hoạt của anh đều ở đây. Ăn, ngủ, học, làm việc\dots\ hết cả trong này. Anh quen kiểu sống như vậy rồi, mấy đứa không cần thấy tội nghiệp cho anh đâu.''

Câu trả lời của cậu khiến một vài người bật lên tiếng cười trừ, nhưng trước khi ai đó kịp nói gì thêm, A-chan đột ngột chen vào: ``Vậy nhà cậu có bao nhiêu người ở vậy?''

Chưa kịp để cậu trả lời, tiếng bước chân vang lên từ cầu thang. Một giọng phụ nữ lớn tuổi vọng lên: ``\vie{Có chuyện gì mà ở trên này rôm rả thế?}''

Cánh cửa mở ra, và người phụ nữ nhà Hikawa - Hikawa Kaien - bước vào cùng chồng, cả hai lập tức tròn mắt trước cảnh tượng đông đúc trước mắt. Bà cất lời, giọng ngạc nhiên: ``\vie{Có tận hai mươi người đang tập trung ở đây à?}''

Ryujin, thợ cơ khí lão luyện nhà Hikawa, cũng không giấu nổi sự bất ngờ: ``\vie{Con\dots\ đây là ai thế?}''

Cậu con trai không hề nao núng, trả lời ngay: ``\vie{À, đồng nghiệp của con ở bên \hll\ đây ạ, bố mẹ.}'' Người bố, nhớ lại cuộc trò chuyện ban nãy, liền gật gù, quay về phía những người vừa tới, nói tiếng Nhật: ``À, cái bên công ty của nó sang làm phim à?''

A-chan ngạc nhiên: ``Sao chú biết?''

``Vừa nãy nó nói với cô chú,'' ông trả lời, giọng điềm đạm nhưng vẫn không giấu được vẻ tò mò khi đảo mắt nhìn quanh phòng. Kaien tiếp tục: ``Nhưng mà sao đông thế? Đây là cả công ty à?''

Kuon bối rối gãi đầu, cười nhạt: ``Nói có lẽ bố mẹ sẽ không tin, nhưng\dots\ đây là tất cả số cô hề-- à không, tất cả số idol mà công ty đang có đấy ạ\dots''

``Kuon-senpai định nói chúng ta là hề à\dots'' Miko tặc lưỡi, rõ ràng không tỏ vẻ vui vẻ gì khi bị cậu nói vậy, nhưng Mel vội ngăn cô lại.

Nàng Thần tượng, từ trong đám đông, giơ tay phát biểu: ``A-nô\dots\ Cháu nghĩ chúng ta nên tìm một chỗ nào đó rộng rãi hơn để trò chuyện\dots''

Nàng Sói đen cũng tiếp lời, khẽ cau mày: ``Cháu đồng ý. Thế này thì mọi người khó thở lắm\dots''

Sau khi hội ý chớp nhoáng, cả nhóm quyết định kéo nhau xuống tầng ba để tiện trò chuyện và\dots\ hít thở thoải mái hơn.

\begin{center}
  \uline{Tầng 3 - nhà Hikawa.}
\end{center}

Cả nhà Kuon, gồm cả Ryuuhou, em gái của cậu, đã về phòng khách và ngồi ở trung tâm. Các thành viên của \hll\ bao quanh thành một vòng tròn, khiến không khí trở nên vô cùng đông vui nhưng cũng có chút trang nghiêm.

``Bây giờ chú muốn biết tất cả các cháu là ai,'' Ryujin nói, ánh mắt ông lướt qua từng khuôn mặt. Nàng Thần tượng liền giơ tay đứng dậy, bắt đầu giới thiệu bản thân mình: ``Cháu là Tokino Sora, thành viên Thế hệ thứ 0 của \hll\ ạ.''

Kuon ngay lập tức chen vào: ``Hay còn được biết đến là Thánh thần.'' Cô cười gượng, khẽ xua tay: ``E-Em không toàn năng đến mức thế đâu\dots''

Kaien cười hiền, chỉ tay về phía Haato: ``Còn cháu kia thì sao? Cháu tên gì?'' Cô nàng Song tính nhanh chóng đứng lên, hăng hái: ``Cháu là Akai Haato, nhưng mọi người cũng hay gọi cháu là Haachama ạ!''

Lần lượt từng người trong nhóm đều tự giới thiệu bản thân. Sau mỗi lời giới thiệu, cậu Tổng không quên thêm vào biệt danh mà cậu hay gọi họ ở công ty. Ai vui vẻ chấp nhận như Haato hay Pekora thì cười toe toét, còn những người không thích lắm như Roboco (bị gọi là ``PON''), Aqua (bị gọi là ``Baqua''), hay Fubuki (bị gọi là ``Cáo Mèo'') thì chỉ biết lườm Kuon nhưng không phản đối mạnh mẽ, vì dù sao đó cũng là một cái biệt danh rất vô hại để đặt cho dễ nhớ và dễ phân biệt mọi người với nhau.

Khi phần giới thiệu đã xong, Ryuuhou quay sang nàng ``Tăng ca'', ánh mắt đầy tò mò: ``Vậy mọi người đến đây để làm gì?''

Cô mỉm cười đáp: ``Các chị tới đây để quay phim về văn hóa của nhà Hikawa-san, em ạ.''

Ryuuhou quay đầu nhìn anh trai với vẻ khó hiểu; ``Thế là thế nào ạ?'' Cậu giải thích: ``Nói ngắn gọn là họ muốn ghi lại cảnh ăn Tết nhà mình. Nhưng anh nghĩ quay thì ít, mà ăn chơi thì nhiều thôi.''

Bình luận của cậu khiến cả phòng bật cười khúc khích. Nhưng không phải ai cũng cảm thấy đồng tình với ý tưởng của cậu.

Matsuri liền phản đối: ``Sao anh lại nghĩ như thế chứ?'' Miko thì nhìn cậu  đầy nghi hoặc: ``Anh có vẻ không tin tưởng bọn này nhỉ?''

Cậu con trai nhún vai: ``Có bao giờ mọi người làm gì đó một cách nghiêm túc đâu.'' Nàng Sói đen, dù cảm thấy không thoải mái gì với lời bình của cậu, cũng chỉ chậm rãi gật đầu: ``Anh ấy nói cũng đúng mà. Hồi nãy chúng ta làm phim, có khi nào suy nghĩ thấu đáo đâu\dots''

Câu nói của Mio khiến cả nhóm ngượng ngùng, không ai cãi lại được. Bố mẹ Kuon thì cười khẽ, còn cậu chỉ biết thở dài.

Cô em gái của Kuon nghiêng đầu hỏi nàng Thần tượng: ``Như Onii-san nói, mọi người đến đây để ghi lại cảnh ăn Tết nhà bọn em ạ?''

Sora mỉm cười, nhẹ nhàng đáp: ``Đúng thế! Ngoài ra đây cũng là dịp để chúng em hiểu hơn về văn hóa địa phương, từ đó có thể kết nối tốt hơn với các kaigai-niki (\ruby{海}{かい}\ruby{外}{がい}ニキ, fan nước ngoài).''

Người mẹ gật đầu, ánh mắt ánh lên vẻ thấu hiểu. Bà nhìn quanh tất cả mọi người, rồi cất giọng hỏi: ``Vậy các cháu có biết `Tết' là gì không?''

Các cô gái bắt đầu chia sẻ những gì mình biết. Chủ yếu là các phong tục truyền thống, những món ăn đặc trưng và ý nghĩa của ngày Tết.

Người bố sau khi nghe mọi người trình bày hiểu biết của mình về ngày lễ đặc biệt đó, liền bật cười: ``Các cháu tìm hiểu cũng kỹ đấy! Nhưng mỗi vùng miền lại có phong tục riêng, nên các cháu còn nhiều điều để học lắm.''

Cậu Tổng gật đầu đồng tình: ``Mọi người nói khá chi tiết rồi. Nhưng ở cả hai miền thì vẫn có những nét đặc trưng riêng.'' Cậu quay sang Ayame, tò mò hỏi: ``Ayame-chan, em từng đón Tết ở đây chưa?''

Nàng Quỷ sừng gật đầu: ``Em từng đến khoảng 3-4 năm trước, trước khi vào \hll-da zo! Hồi đó em ở tỉnh \ruby{Nishimitsugi}{Sài Gòn} cơ!''

Cậu ngạc nhiên: ``Vậy em từng đón Tết ở miền Nam rồi à?'' Ayame cười khẽ: ``Có lẽ vậy\dots? Nhưng đây là lần đầu em đến Kawachi. Mọi thứ ở đây khác hẳn luôn!''

Noel tò mò hỏi: ``Miền Nam? Hình như Sora-chan với A-chan cũng vừa nói với chúng em như vậy.''

Kuon tiếp tục giải thích: ``Chúng ta đang ở miền Bắc của Etsunan. Còn miền Trung cũng có Tết, nhưng phong tục ở đó chia thành hai nửa, mỗi vùng một khác.''

Cả phòng đồng thanh *ồ* lên, cảm thấy hứng thú với những điều vừa được chia sẻ. Nàng Hiệp sĩ khẽ nói: ``Giờ em mới biết Tết cũng được ăn ở cả miền Trung nữa đấy\dots''

Ryujin mỉm cười, gật đầu. Ông chậm rãi hỏi: ``Vậy các cháu có biết ý nghĩa thật sự của ngày Tết là gì không?''

Miko mạnh dạn giơ tay phát biểu, dù giọng điệu hơi ngập ngừng: ``Có phải là thời điểm báo hiệu mùa xuân tới không ạ, ne?''

Cô em gái của cậu ghé sát vào tai của Kuon: ``\hl{Ne?}'' Cậu đáp: ``\hl{Thói quen của em ấy. Đừng hỏi.}''

Người thợ cơ khí gật đầu: ``Ý của cháu đúng rồi đấy. Còn gì nữa không?'' Nàng Ma cà rồng tiếp lời: ``Là thời khắc chuyển giao giữa năm cũ và năm mới ạ?''

Người đàn bà nội trợ mỉm cười hiền hậu: ``Không sai.''

Pekora, vẻ mặt đầy hào hứng, bổ sung: ``Là chuẩn bị cho mùa gieo gặt mới, peko?'' Cậu con trai quay sang nhìn cô nàng Thỏ, có vẻ bất ngờ: ``Anh không nghĩ là em sẽ trả lời đúng đấy.''

``Quê em trồng cà rốt, nên em đoán thế, peko,'' nàng Thỏ đáp một cách hóm hỉnh, vênh mặt lên.

Tiếp đó, nàng tóc bay Aki lên tiếng: ``Cháu nghĩ đó là dịp để chúng ta nhớ về cội nguồn, phải không ạ?'' Câu nói này khiến Matsuri nhíu mày bối rối. Cô nàng tiếp lời một cách đầy triết lý: ``Cội nguồn như tế bào ban đầu tạo ra con người, dần dần tiến hóa thành xã hội ngày nay\dots''

Mio vội vàng lên tiếng: ``Em không nghĩ ý nghĩa của Tết lại sâu xa như thế đâu, Matsuri-chan\dots''

Cậu con trai cũng lắc đầu, thở dài: ``Ờm\dots\ `Nhớ về cội nguồn' là đúng, nhưng không cần phải phân tích kiểu `đại khai trí tuệ' như vậy đâu.''

Người mẹ Kaien mỉm cười giải thích: ``Cội nguồn ở đây là quê hương của chúng ta, nơi các cháu sinh ra và lớn lên, nơi luôn giữ trong mình những kỷ niệm về gia đình và tuổi thơ.'' Kuon liền chen vào: ``Mẹ ơi, mẹ nói thế họ khó hiểu lắm. Nói ngắn gọn hơn đi mẹ.''

Bà bật cười: ``Cô tưởng là các cháu sẽ hiểu chứ?'' Cậu con trai quay sang nhóm idol, khẽ nhướng mày: ``Mọi người có hiểu mẹ anh nói không?''

Aki lập tức giơ tay, ánh mắt sáng rực: ``Hiểu chứ! Là biển cả!'' Haato thở dài, vội sửa lại: ``Không phải, là quê nhà của chúng ta đó, Haachama-desu!''

Nàng Dị giới bối rối: ``Quê nhà?'' Ryuuhou cười khúc khích, chêm vào: ``Là nơi mà bố mẹ sinh ra chúng ta đó, Aki-senpai.'' Lúc này, cô mới trố mắt: ``Ồ\dots\ Giờ chị hiểu rồi.''

Cậu con trai mỉm cười: ``Như Haato-chan nói đấy. Tết là dịp để mọi người về thăm quê nhà, đoàn tụ cùng gia đình.''

Khi mọi người đã hiểu được giá trị thật sự của ngày Tết, tất cả hai mươi người con gái đều gật gù. Nàng Thần tượng gốp lời, tỏ vẻ hơi ái ngại: ``Thật ra bọn em làm việc quanh năm nên hiếm khi có dịp về thăm gia đình lắm anh.''

Kuon nghe lời giải thích của cô, khẽ gật đầu, ánh mắt hiện lên sự đồng cảm: ``Anh hiểu. Văn hóa làm việc của công ty bọn em khá bận rộn mà.''

Rushia cất giọng nhỏ nhẹ: ``Nghe có vẻ đầy lễ nghĩa đây\dots'' Kaien mỉm cười động viên mọi người: ``Các cháu cứ trải nghiệm dần là sẽ hiểu thôi.''

A-chan cười khẽ, quay sang cả nhóm và nói: ``Đúng như cô nói ạ. Chính vì thế chúng cháu mới tới đây, để lấy tư liệu phim!''

Ryujin mỉm cười, tiếp tục giải thích: ``Ngày Tết là dịp mọi người tưởng nhớ tổ tiên, tri ân những người đã khuất và sum vầy bên gia đình. Đây còn là thời điểm chào đón mùa xuân, cầu mong một năm mới nhiều sức khỏe, may mắn và hạnh phúc.''

Người vợ của ông nói thêm: ``Tết còn là thời điểm để bỏ lại mọi điều không vui của năm cũ, hướng tới tương lai với những khởi đầu mới.''

Cô em gái của Kuon cũng lên tiếng: ``Tết còn là khoảng thời gian chúng ta xả hơi sau một năm làm việc vất vả nữa, các chị à!''

Cả nhóm chăm chú lắng nghe, thỉnh thoảng gật đầu khi hiểu thêm ý nghĩa của Tết. Không khí trong phòng ấm cúng và đầy hứng thú.

\ct

Sau một hồi nghe gia đình nhà Hikawa giải thích về ý nghĩa của ngày Tết, tất cả mọi người cũng đã hiểu rõ hơn về ngày lễ đặc biệt thường niên tại xứ người.

``Tất nhiên, trăm nghe không bằng một thấy,'' Kaien mỉm cười, ``nên cô nghĩ rằng có rất nhiều hoạt động ngày Tết các cháu có thể tham gia. Bắt đầu với\dots''

Cả hai mươi cô gái nín thở, ánh mắt háo hức hướng về người mẹ. Nhưng chỉ trong chốc lát, họ đồng loạt giật mình khi cả nhà Hikawa đồng thanh hô: ``\textbf{Dọn dẹp nhà cửa!}''

Không ai thực sự ngạc nhiên, nhất là Mio và A-chan. Nàng Đeo kính điềm nhiên nhận xét: ``Cái này giống như khi chúng ta dọn đống hỗn độn mà mọi người gây ra sau mỗi tập \textit{hologra} thôi mà.''

Nàng Sói đen than thở, khẽ rùng mình khi cô nhớ lại những lần mà các thành viên ở \hll\ phá tanh bành cả văn phòng, để rồi phải để chính cô, A-chan (và đôi khi là Kuon) dọn dẹp. ``Dọn dẹp đến mức liệt giường luôn\dots''

Người bạn thân của cô tỏ vẻ lo lắng: ``Liệu chúng ta có phải làm những công việc như thế không nhỉ?''

Không để cho các cô gái đoán già đoán non, cậu Tổng thẳng thừng đáp ngay: ``Câu trả lời là có. Và chúng ta sẽ làm nhiều hơn thế.''

Tất cả mọi người đều đổ mồ hôi hột. Subaru hỏi, tỏ vẻ tò mò, dù vẫn còn hơi lo lắng: ``Cụ thể là như thế nào?'' Cậu không ngần ngại đáp: ``Tất cả mọi thứ. Những gì mà mọi người thực hành ngay bây giờ. Có cả một núi việc đấy.''

Thấy mọi người quay sang mình như chờ lời giải thích, Ryujin chậm rãi nói: ``Để đón Tết, việc đầu tiên các cháu cần làm là dọn dẹp nhà cửa. Không ai muốn đón năm mới trong một căn nhà bừa bộn, đúng không?''

Cả Shion và Matsuri đều không tỏ vẻ thoải mái gì khi họ buộc phải động tay động chân làm. Nàng Phù thủy nhíu mày: ``Không thể dời sang mấy ngày sau được sao?'', còn nàng Cổ động phụ họa đầy e dè: ``Ừm\dots\ Chị cảm thấy chúng ta sẽ phải làm nhiều việc đó\dots''

Ngay lập tức, Kuon phản đối: ``Không được. Ngày mai là 30 Tết rồi, mà chúng ta còn cả đống việc cần làm. Đợi đến ngày mai ngày kia thì muộn mất.''

Cô nàng Lễ hội nhanh nhảu gợi ý: ``Vậy mùng 1 thì sao?'' Câu hỏi này lại liền bị người bố của Kuon phản bác: ``Ai lại đi dọn nhà vào mùng 1 Tết bao giờ!''

``Tại sao thế ạ?'' hai người họ bĩu môi, tỏ vẻ mệt mỏi.

Người mẹ của Kuon mỉm cười giải thích: ``Ở đây bọn cô kiêng dọn nhà vào mùng 1 Tết, cháu à. Người ta tin rằng dọn dẹp vào ngày đó chẳng khác nào quét hết may mắn ra khỏi nhà.''

Shion vẫn tỏ vẻ không phục: ``Nghe có vẻ\dots\ ngớ ngẩn thế nào ấy\dots'' Thấy vậy, Kaien tiếp tục nhẹ nhàng khuyên nhủ: ``Việc hôm nay chớ để ngày mai. Làm xong sớm thì nghỉ sớm, đúng không?''

Nàng Vu nữ gật đầu, tỏ vẻ nghiêm túc: ``Cô ấy nói đúng đấy, ne. Kami-sama sẽ không vui nếu chúng ta đùn công việc sang ngày khác đâu, ne.''

Nàng Sói đen cũng đồng quan điểm với cô ấy, trấn an mọi người: ``Với hai mươi người thế này, tôi nghĩ chúng ta sẽ làm xong sớm thôi.''

Những thành viên khác đồng loạt gật đầu đồng ý. Shion và Matsuri, dù không mấy vui vẻ, cũng chỉ có thể ngậm ngùi tuân theo.

Thế là, công cuộc dọn dẹp nhà cửa đón Tết của gia đình Hikawa chính thức bắt đầu. Không khí hối hả, bận rộn nhưng cũng đầy tiếng cười và sự phấn khởi chuẩn bị cho ngày lễ truyền thống quan trọng nhất trong năm.

``\textit{Đâu ai ngờ rằng việc đầu tiên để có một ngày Tết ấm cúng lại\dots\ vất vả như vậy chứ?}''

\end{document}